\chapter{初始轨迹构造}

\section{初始化必要性}
RSSI 观测存在随机跳变，直接在全局平面上优化容易出现局部极小。合理初值能显著降低迭代次数并提升鲁棒性，因此本文采用分阶段初始化。

\section{单帧粗定位}
在时刻 $t$，选择信号最强的 $K$ 个信标（通常 $K=3\sim5$），最小化
\[
\min_{\bm p_t}\sum_{i=1}^{K}\left(d_i(\bm p_t)-\hat d_i\right)^2.
\]
该问题通过网格搜索得到粗解 $\tilde{\bm p}_t$，作为后续窗口优化起点。

\section{轨迹级初始化}
将单帧粗解串联为初始轨迹 $\{\tilde{\bm p}_t\}$，再执行两步修正：
\begin{enumerate}
  \item 中值滤波去除突发跳点；
  \item 根据速度上界裁剪异常位移段。
\end{enumerate}
得到平滑初始轨迹 $\bm X^{(0)}$。

\section{可行域投影初始化}
若某时刻初值落入障碍区，则通过投影算子
\[
\bm p_t^{(0)}=\Pi_{\Omega}(\tilde{\bm p}_t)
\]
映射到最近可行点。该步骤能避免后续优化初期即受到硬约束冲突影响。

\section{缺失观测时的补点策略}
当时刻 $t$ 可用信标数量不足（例如少于 2 个）时，采用短时外推：
\[
\bm p_t^{(0)}=\bm p_{t-1}^{(0)}+\alpha(\bm p_{t-1}^{(0)}-\bm p_{t-2}^{(0)}),
\]
其中 $\alpha\in[0,1]$ 为衰减系数。该策略可保持轨迹连续。

\section{初始化质量评估}
定义初值质量指标：
\[
E_0=\frac{1}{T}\sum_{t=1}^{T}\min_i\|\bm p_t^{(0)}-\bm b_i\|_2,
\]
并结合速度一致性指标筛选初值参数。实验表明，初始化阶段处理充分时，主优化迭代次数可降低约 30\%。
